% C4 Architecture Visualization Documentation Suite | original engineering appendix
% File: reference/diagram-review-checklist-and-acceptance-criteria.tex
% Standalone pdfLaTeX root; no shell escape or external style file.
\documentclass[10pt,letterpaper]{article}
\usepackage[T1]{fontenc}
\usepackage[utf8]{inputenc}
\usepackage{lmodern,microtype}
\usepackage[letterpaper,margin=0.82in,headheight=15pt,headsep=18pt,footskip=26pt]{geometry}
\usepackage[table]{xcolor}
\usepackage{booktabs,tabularx,array,longtable,enumitem}
\usepackage{fancyhdr,lastpage,titlesec}
\usepackage{hyperref,xurl,amssymb,amsmath}
\definecolor{C4Navy}{HTML}{17364B}
\definecolor{C4Teal}{HTML}{137C85}
\definecolor{C4Gray}{HTML}{526775}
\definecolor{C4Light}{HTML}{EEF4F5}
\definecolor{C4Line}{HTML}{CAD8DE}
\color{C4Navy}
\hypersetup{colorlinks=true,linkcolor=C4Teal,urlcolor=C4Teal,
pdftitle={Architecture Diagram Review Checklist and Measurable Acceptance Criteria},
pdfauthor={Architecture Engineering Documentation Suite}}
\setlength{\parindent}{0pt}
\setlength{\parskip}{5pt}
\setlength{\emergencystretch}{2em}
\raggedbottom
\setlist[itemize]{leftmargin=1.5em,topsep=3pt,itemsep=3pt}
\setlist[enumerate]{leftmargin=1.6em,topsep=3pt,itemsep=3pt}
\titleformat{\section}{\Large\bfseries\sffamily\color{C4Navy}}{\thesection}{0.6em}{}
\titleformat{\subsection}{\large\bfseries\sffamily\color{C4Teal}}{\thesubsection}{0.6em}{}
\titlespacing*{\section}{0pt}{16pt}{7pt}
\titlespacing*{\subsection}{0pt}{12pt}{5pt}
\pagestyle{fancy}
\fancyhf{}
\fancyhead[L]{\scriptsize\sffamily C4 / ARCHITECTURE VISUALIZATION}
\fancyhead[R]{\scriptsize\sffamily REVIEW STANDARD / APPENDIX}
\fancyfoot[L]{\scriptsize\color{C4Gray}ORIGINAL ENGINEERING SUPPLEMENT / 2026}
\fancyfoot[R]{\scriptsize\thepage\ / \pageref*{LastPage}}
\renewcommand{\headrulewidth}{0.4pt}
\newcolumntype{Y}{>{\raggedright\arraybackslash}X}
\newcommand{\Gate}{\textsuperscript{\(\blacklozenge\)}}
\newcommand{\CheckRow}[3]{\textbf{\textsf{#1}} & #2 & #3\\[3pt]}
\newcommand{\ChecklistHeader}{\toprule
\textbf{ID} & \textbf{Question} & \textbf{Observable acceptance criterion}\\\midrule
\endfirsthead
\toprule\textbf{ID} & \textbf{Question (continued)} & \textbf{Acceptance criterion}\\\midrule
\endhead
\midrule\multicolumn{3}{r}{\footnotesize Continued on next page}\\\endfoot
\bottomrule\endlastfoot}
\begin{document}
\begin{titlepage}
\renewcommand{\thepage}{title-\arabic{page}}
\vspace*{0.65in}
{\sffamily\bfseries\small\color{C4Teal}SOFTWARE ARCHITECTURE / CONTROLLED APPENDIX}\par\vspace{0.23in}
{\sffamily\bfseries\fontsize{28}{32}\selectfont Architecture Diagram Review Checklist}\par\vspace{0.18in}
{\sffamily\fontsize{17}{22}\selectfont Measurable Acceptance Criteria, Evidence Protocol and Quality Gates}\par\vspace{0.33in}
\textcolor{C4Teal}{\rule{\linewidth}{1.2pt}}\par\vspace{0.26in}
\begin{tabularx}{\linewidth}{@{}>{\raggedright\arraybackslash}p{0.25\linewidth}Y@{}}
\textbf{Document ID} & ARCH-C4-APP-A-DIAGRAM-REVIEW\\
\textbf{Edition / status} & 1.0 / proposed reusable review standard\\
\textbf{Updated} & 16 November 2026\\
\textbf{Scope} & C4 context, container, component, code, dynamic, deployment and landscape; adaptable to other architecture diagrams\\
\textbf{Audience} & Architects, developers, AppSec, SRE, reviewers and approvers\\
\end{tabularx}
\vspace{0.35in}
\fcolorbox{C4Line}{C4Light}{\parbox{0.93\linewidth}{\footnotesize
\textbf{Source boundary.} The supplied 21-question taxonomy (General, Elements, Relationships) is adapted and cross-checked against Simon Brown's official C4 guidance. The further 42 criteria, measurements, gates, and workflow are independently authored engineering supplements, not normative C4 or ISO/IEC/IEEE 42010 requirements.}}
\vfill
{\small\color{C4Gray}Passing this proposed standard documents communication quality, not runtime correctness, security certification or formal standards conformance.}\par
\end{titlepage}
\tableofcontents
\clearpage
\appendix
\renewcommand{\thetable}{A-\arabic{table}}
\section{Diagram Review Checklist and Measurable Acceptance Criteria}
\label{app:diagram-review}
\subsection{Purpose and application}
This appendix translates questions about diagram interpretation into verifiable tests for design review, architectural governance, onboarding, threat analysis and publishing. Tables A-1 through A-3 preserve the original three-category organization. Further tables extend review to C4 abstraction, cross-view fidelity, trust boundaries, runtime, accessibility and evidence governance.

\textbf{Freeze the artifact before review.} Each record must identify the diagram ID, author, independent reviewer, source revision, rendered artifact, audience, viewpoint, scope, environment or target state and review date. Do not treat a proposed architecture as verified production state.

\begin{itemize}
\item \textbf{Pass:} the acceptance statement is supported by an explicit check or cited artifact.
\item \textbf{Fail:} missing, inaccurate, contradictory or unverifiable content; unassessed applicable tests count as Fail.
\item \textbf{N/A:} the property genuinely does not apply, with a written rationale confirmed by the reviewer; never a substitute for missing evidence.
\item \textbf{Critical gate:} a diamond (\(\blacklozenge\)) means the applicable test must pass regardless of average score.
\end{itemize}
The quantifier 100\% means all in-scope instances, not an arbitrary sample. When there are zero in-scope items, record N/A and the population rationale. All thresholds below are proposed local engineering requirements unless explicitly identified as an existing suite scorecard rule.

\subsection{General}
\begin{longtable}{@{}>{\raggedright\arraybackslash}p{0.09\linewidth}>{\raggedright\arraybackslash}p{0.42\linewidth}>{\raggedright\arraybackslash}p{0.42\linewidth}@{}}
\caption{General (4 criteria)}\label{tab:review-G}\\
\ChecklistHeader
\CheckRow{G-01}{Does the diagram have a meaningful title?}{Title names the view type and system/domain and matches the registered artifact.}
\CheckRow{G-02\Gate}{Can a reviewer identify the view type?}{The declared type is visible and agrees with the abstraction shown.}
\CheckRow{G-03\Gate}{Is diagram scope unambiguous?}{System of interest, boundary, environment/current-versus-target state and material exclusions are recorded.}
\CheckRow{G-04}{Does it provide a notation key?}{Every non-obvious color, shape, icon, border, edge and arrowhead is defined.}
\end{longtable}

\subsection{Elements}
\begin{longtable}{@{}>{\raggedright\arraybackslash}p{0.09\linewidth}>{\raggedright\arraybackslash}p{0.42\linewidth}>{\raggedright\arraybackslash}p{0.42\linewidth}@{}}
\caption{Elements (10 criteria)}\label{tab:review-E}\\
\ChecklistHeader
\CheckRow{E-01}{Does every element have a name?}{100\% of elements have legible names; duplicates use stable IDs for disambiguation.}
\CheckRow{E-02\Gate}{Is the type of every element clear?}{100\% have correct declared abstraction types; intentional mixed-level supporting items are distinguished.}
\CheckRow{E-03}{Can a reader explain what each element does?}{Each material element has a concise responsibility statement and no unexplained generic box remains.}
\CheckRow{E-04}{Are applicable technology choices identified?}{Every material container/component identifies technology or is explicitly marked unknown/proposed.}
\CheckRow{E-05}{Are acronyms, abbreviations, and code names understood?}{Every audience-specific term is expanded on first use or in the legend/glossary.}
\CheckRow{E-06}{Do all element colors have defined meanings?}{Each semantic color has a documented, consistent meaning or is designated decorative.}
\CheckRow{E-07}{Do all shapes have defined meanings?}{Each semantic shape maps to a documented type, with an additional discriminator if needed.}
\CheckRow{E-08}{Are all icons understandable?}{Every semantic icon is labeled or explained; icon recognition is not the only cue.}
\CheckRow{E-09}{Are all border styles understandable?}{Every semantic solid, dotted, or dashed boundary is defined and used consistently.}
\CheckRow{E-10}{Are element sizes understandable?}{Size encoding is explained with its measure or expressly identified as decorative.}
\end{longtable}

\subsection{Relationships}
\begin{longtable}{@{}>{\raggedright\arraybackslash}p{0.09\linewidth}>{\raggedright\arraybackslash}p{0.42\linewidth}>{\raggedright\arraybackslash}p{0.42\linewidth}@{}}
\caption{Relationships (7 criteria)}\label{tab:review-R}\\
\ChecklistHeader
\CheckRow{R-01\Gate}{Does every arrow describe relationship intent?}{100\% of edges have specific verb-based labels, not only generic labels such as Uses.}
\CheckRow{R-02\Gate}{Does every label match arrow direction?}{100\% of origins/targets agree with text; two-way interactions have two explained directions.}
\CheckRow{R-03\Gate}{Are relationship technologies/protocols identified where applicable?}{Every interprocess container edge states material transport/protocol or an explicit unknown.}
\CheckRow{R-04}{Are all relationship acronyms explained?}{All audience-specific edge abbreviations appear in the legend or glossary.}
\CheckRow{R-05}{Are relationship colors explained?}{Every semantic edge color is defined and still identifiable without color.}
\CheckRow{R-06}{Are arrowhead conventions explained?}{Arrowheads represent the intended direction and semantics consistently.}
\CheckRow{R-07}{Are line styles explained?}{Solid/dashed/dotted edge meanings are defined and applied consistently.}
\end{longtable}

\subsection{Abstraction and cross-view correspondence}
\begin{longtable}{@{}>{\raggedright\arraybackslash}p{0.09\linewidth}>{\raggedright\arraybackslash}p{0.42\linewidth}>{\raggedright\arraybackslash}p{0.42\linewidth}@{}}
\caption{Abstraction and cross-view correspondence (8 criteria)}\label{tab:review-X}\\
\ChecklistHeader
\CheckRow{X-01\Gate}{Does the view honor the declared abstraction level?}{No in-scope node violates the selected zoom; supporting exceptions are explicitly labeled.}
\CheckRow{X-02\Gate}{Are ownership and containment boundaries correct?}{Each component belongs to its container; each container to its system; external actors remain outside.}
\CheckRow{X-03}{Do shared entities retain stable identities?}{100\% of repeated nodes use the same stable ID/type across linked views or a recorded rename mapping.}
\CheckRow{X-04}{Do adjacent views reconcile?}{Material context edges map to container interactions; component edges map to owning interfaces or documented exceptions.}
\CheckRow{X-05}{Are significant omissions disclosed?}{Every known material simplification is recorded with rationale, owner and impact.}
\CheckRow{X-06}{Are system, deployment, network and trust boundaries distinct?}{Each boundary type is visibly labeled; nesting alone is not treated as proof of trust.}
\CheckRow{X-07}{Can material components be mapped to code or catalog records?}{Current-state material components have a source package, repository, or service-catalog reference.}
\CheckRow{X-08}{Are instance-count and redundancy claims verified?}{Multiplicity, quorum and replicas have dated deployment evidence or are clearly target-state.}
\end{longtable}

\subsection{Security, privacy and data assurance}
\begin{longtable}{@{}>{\raggedright\arraybackslash}p{0.09\linewidth}>{\raggedright\arraybackslash}p{0.42\linewidth}>{\raggedright\arraybackslash}p{0.42\linewidth}@{}}
\caption{Security, privacy and data assurance (8 criteria)}\label{tab:review-S}\\
\ChecklistHeader
\CheckRow{S-01\Gate}{Are material trust crossings represented?}{All material ingress, egress and zone-crossing flows appear or link to a named threat-model/DFD artifact.}
\CheckRow{S-02\Gate}{Are sensitive data flows classified?}{Every material sensitive, regulated or privileged flow has an explicit classification in a linked register.}
\CheckRow{S-03\Gate}{Are security-critical assumptions and owners traceable?}{Each material trust crossing names control assumptions and an accountable owner, or has an unresolved finding.}
\CheckRow{S-04}{Are public/partner interfaces inventoried?}{100\% of material external entry points map to an API, event or integration inventory.}
\CheckRow{S-05}{Are data-store duties identifiable?}{Each material data store has an identified data category, owner and applicable retention/deletion reference.}
\CheckRow{S-06}{Are relevant abuse and failure paths reviewable?}{High-risk interactions show or link at least one relevant denial, timeout, replay or validation failure path.}
\CheckRow{S-07}{Is sensitive information excluded from published output?}{Inspection finds zero credentials, tokens, private records or unauthorized internal-only details.}
\CheckRow{S-08}{Are unresolved security assertions tracked?}{Each unresolved assertion has a finding ID, owner, severity and due date.}
\end{longtable}

\subsection{Operational and runtime fidelity}
\begin{longtable}{@{}>{\raggedright\arraybackslash}p{0.09\linewidth}>{\raggedright\arraybackslash}p{0.42\linewidth}>{\raggedright\arraybackslash}p{0.42\linewidth}@{}}
\caption{Operational and runtime fidelity (6 criteria)}\label{tab:review-O}\\
\ChecklistHeader
\CheckRow{O-01}{Are observed, inferred and proposed assertions separated?}{Each runtime/deployment claim declares status; current-state claims link dated evidence.}
\CheckRow{O-02}{Are critical dependencies accounted for?}{Every critical service-catalog or impact-assessment dependency is shown or cross-referenced.}
\CheckRow{O-03}{Are asynchronous relationships accurate?}{Material topics/queues identify producer, consumer, direction and relevant delivery assumptions.}
\CheckRow{O-04}{Are reliability mechanisms substantiated?}{Retry, failover, replication and timeout claims point to configuration or tests.}
\CheckRow{O-05}{Is operational ownership traceable?}{Critical runtime nodes have owning teams and monitoring/runbook references when operations are in scope.}
\CheckRow{O-06}{Are deployment environments explicit?}{The environment, region/zone and material differences are named; environments are not silently conflated.}
\end{longtable}

\subsection{Accessibility and visual comprehension}
\begin{longtable}{@{}>{\raggedright\arraybackslash}p{0.09\linewidth}>{\raggedright\arraybackslash}p{0.42\linewidth}>{\raggedright\arraybackslash}p{0.42\linewidth}@{}}
\caption{Accessibility and visual comprehension (5 criteria)}\label{tab:review-A}\\
\ChecklistHeader
\CheckRow{A-01}{Does a monochrome export preserve meaning?}{All semantic distinctions remain recognizable without color.}
\CheckRow{A-02}{Can labels be read at publication scale?}{Rendered labels are legible at intended 100\% viewing size, with no cropping or overlap.}
\CheckRow{A-03}{Is the reading order understandable?}{An independent intended reader can identify the entry point and principal interactions without oral help.}
\CheckRow{A-04}{Are crossed edges and dense clusters unambiguous?}{Critical paths have no confusing intersections, or connector anchors/split views remove ambiguity.}
\CheckRow{A-05}{Is an accessible text alternative available where needed?}{Published packet includes a narrative or tabular element/relationship inventory.}
\end{longtable}

\subsection{Evidence, versioning and review governance}
\begin{longtable}{@{}>{\raggedright\arraybackslash}p{0.09\linewidth}>{\raggedright\arraybackslash}p{0.42\linewidth}>{\raggedright\arraybackslash}p{0.42\linewidth}@{}}
\caption{Evidence, versioning and review governance (8 criteria)}\label{tab:review-V}\\
\ChecklistHeader
\CheckRow{V-01\Gate}{Are claims bound to a specific version?}{Diagram ID, immutable source revision, render reference, author and review date are recorded.}
\CheckRow{V-02}{Can the diagram be regenerated?}{Source file, toolchain and build instructions reproduce the expected output.}
\CheckRow{V-03\Gate}{Was the diagram independently reviewed?}{A reviewer other than the author records a date, disposition and critical-gate check.}
\CheckRow{V-04}{Are evidence sources sufficiently current?}{Each source carries an observation date and the reviewer approves its currency for the change rate.}
\CheckRow{V-05}{Are model/code/deployment discrepancies tracked?}{Every known material discrepancy has a stable issue, owner and disposition.}
\CheckRow{V-06}{Are reader roles and concerns identified?}{At least one audience role and the questions the diagram answers are recorded.}
\CheckRow{V-07}{Are re-review triggers defined?}{Boundary, interface, protocol, trust flow, critical dependency and deployment changes trigger review.}
\CheckRow{V-08}{Has distribution clearance been checked?}{Licensing, confidentiality, naming, and evidence-sharing constraints are assessed before publication.}
\end{longtable}

\subsection{Conditional view-type probes}
Apply the probe matching each declared view; for a non-C4 diagram, state the analogous viewpoint-specific test rather than marking all seven as required.
\begin{longtable}{@{}>{\raggedright\arraybackslash}p{0.11\linewidth}>{\raggedright\arraybackslash}p{0.24\linewidth}>{\raggedright\arraybackslash}p{0.57\linewidth}@{}}
\caption{Diagram-type acceptance tests (seven conditional criteria)}\label{tab:review-C}\\
\toprule\textbf{ID} & \textbf{View} & \textbf{Acceptance test}\\\midrule\endfirsthead
\toprule\textbf{ID} & \textbf{View} & \textbf{Acceptance test}\\\midrule\endhead
\bottomrule\endlastfoot
\textbf{\textsf{C-01}} & System context & Exactly one system of interest is designated; primary users and directly connected external systems have labeled roles and interactions.\\[3pt]
\textbf{\textsf{C-02}} & Container & Boundary encloses the main applications/data stores; responsibilities, technologies, and interprocess protocols are explicit.\\[3pt]
\textbf{\textsf{C-03}} & Component & Exactly one primary container is decomposed; important components, interfaces and external collaborators map to that container.\\[3pt]
\textbf{\textsf{C-04}} & Code & Code symbols match a named revision and programming-language semantics; only details needed by the question are included.\\[3pt]
\textbf{\textsf{C-05}} & Dynamic & Named scenario, initial/final conditions, interaction ordering, initiators and material failure paths agree with static contracts.\\[3pt]
\textbf{\textsf{C-06}} & Deployment & Environment and region/zone are named; nodes, container instances, ingress and boundaries map to logical containers.\\[3pt]
\textbf{\textsf{C-07}} & Landscape & Domain boundary, software systems, major consumers, owner references and cross-system edges match portfolio records.\\[3pt]
\end{longtable}

\subsection{Acceptance metrics and release decisions}
\label{sec:review-acceptance}
There are \textbf{63 tests}: 21 original-category tests, 35 additional cross-cutting checks, and seven conditional view-type probes. Most packets apply exactly one view-type probe, so 63 is \emph{not} the default scoring denominator. Each applicable test requires its ID, result, evidence source, reviewer and finding reference if failed.

\subsubsection{Checklist completeness and coverage}
Let \(N_P\) represent passed applicable tests and \(N_F\) represent failed applicable tests. A justified N/A is excluded. Calculate
\[
N_{\mathrm{app}}=N_P+N_F,\qquad
\mathrm{ChecklistCoverage}=100\frac{N_P}{N_{\mathrm{app}}}\%.
\]
If \(N_{\mathrm{app}}=0\), the review is invalid. \textbf{Proposed gate:} at least 95\% applicable Pass results and 100\% of applicable diamond-marked gates Pass. An unassessed applicable item counts as Fail. An N/A rationale must be approved independently.

\subsubsection{Alignment with the existing suite scorecard}
The companion \emph{Diagram Communication and Semantics Scorecard} defines six weighted dimensions. Keep these weights and 0--4 ratings; checklist coverage is a separate metric.
\begin{table}[htbp]
\centering\small
\caption{Compatible quality dimensions and evidence mapping}\label{tab:review-weights}
\begin{tabularx}{\linewidth}{@{}>{\raggedright\arraybackslash}p{0.48\linewidth}>{\raggedright\arraybackslash}p{0.14\linewidth}Y@{}}
\toprule\textbf{Dimension} & \textbf{Weight} & \textbf{Primary test IDs}\\\midrule
Purpose, audience and scope & 20\% & G-01--G-03, V-06, relevant C\\
Abstraction and boundary correctness & 20\% & E-02, X-01--X-08, relevant C\\
Element names and responsibilities & 15\% & E-01, E-03, E-04, X-07\\
Relationships, protocol and direction & 15\% & R-01--R-03, O-02, O-03\\
Notation, legend and accessibility & 10\% & G-04, E-05--E-10, R-04--R-07, A-01--A-05\\
Evidence currency and reviewability & 20\% & S-01--S-08, O-01/O-04--O-06, V-01--V-08\\\midrule
\textbf{Total} & \textbf{100\%} & Mandatory gates override numerical scores.\\
\bottomrule
\end{tabularx}
\end{table}
\[
\mathrm{WeightedScore}=\sum_{i=1}^{6}w_i\frac{r_i}{4},
\qquad r_i\in\{0,1,2,3,4\},\quad \sum_iw_i=100.
\]
\textbf{Suggested evidence anchors:} 0 = absent/contradicted; 1 = assertion with major gaps; 2 = partially documented; 3 = reviewed and substantially complete; 4 = independently reproducible, current and unambiguous.

Retain the suite's prior bands: 85--100 = eligible for approval after blockers clear; 70--84.99 = changes required; below 70 = rework. The new 95\% check and mandatory gates are \textbf{proposed local additions}, not rules dictated by C4.

\subsubsection{Version-bound approval policy}
Mark \textbf{Approved} only when all of the following hold:
\begin{enumerate}
\item Weighted score is at least 85 and applicable checklist coverage is at least 95\%.
\item Every applicable mandatory gate passes, and there are no unresolved Critical or High findings relevant to the diagram's stated purpose.
\item All remaining failures have actionable issues with owners, due dates and documented nonblocking disposition.
\item An independent reviewer records the exact approved source revision, render, date and publication boundary.
\end{enumerate}
Otherwise record \textbf{Changes required} or \textbf{Rework}, according to the defect severity and weighted score. Score alone cannot approve. Critical tests cannot be bypassed by unjustified N/A.

\subsubsection{Worked example (fictional)}
Ratings of \(4,3,4,3,4,3\) in table order yield
\[
20(4/4)+20(3/4)+15(4/4)+15(3/4)+10(4/4)+20(3/4)
=\boxed{86.25}.
\]
With 56 Pass and two Fail among 58 applicable tests, coverage is \(56/58=96.55\%\) (rounded). This is \emph{eligible}, not automatically approved: critical gates must pass, blockers must close, and independent sign-off is still required.

\subsection{Review evidence and severity}
\begin{tabularx}{\linewidth}{@{}>{\raggedright\arraybackslash}p{0.17\linewidth}Y@{}}
\toprule\textbf{Severity} & \textbf{Interpretation and disposition}\\\midrule
Critical & Wrong system/trust boundary, undisclosed sensitive flow or architecture representation creating immediate unacceptable decision risk; blocks approval.\\
High & Wrong material direction/protocol, omitted critical dependency or unsupported significant deployment/security assertion; blocks approval.\\
Medium & Ambiguous responsibility, noncritical missing evidence, cross-view inconsistency; correct and track according to local policy.\\
Low & Wording or cosmetic discrepancy with no material effect on interpretation; track with an owner.\\
\bottomrule
\end{tabularx}

For every failed, mandatory or sampled passing test, record the \textbf{criterion ID, element/edge ID, immutable source revision, evidence URI/path, observation date, reviewer, finding ID, owner and disposition}. The full result register covers \emph{all} applicable criteria even if detailed evidence is sampled for noncritical passes. Never include credentials or personal records in review output.

\subsubsection{Blank review register}
\begin{longtable}{@{}>{\raggedright\arraybackslash}p{0.13\linewidth}>{\raggedright\arraybackslash}p{0.12\linewidth}>{\raggedright\arraybackslash}p{0.23\linewidth}>{\raggedright\arraybackslash}p{0.18\linewidth}>{\raggedright\arraybackslash}p{0.22\linewidth}@{}}
\caption{Review-result register (copy rows as required)}\label{tab:review-register}\\
\toprule\textbf{ID} & \textbf{Pass/Fail/N/A} & \textbf{Evidence/version} & \textbf{Finding/owner} & \textbf{Reviewer/date}\\\midrule\endfirsthead
\toprule\textbf{ID} & \textbf{Result} & \textbf{Evidence/version} & \textbf{Finding/owner} & \textbf{Reviewer/date}\\\midrule\endhead
\bottomrule\endlastfoot
\rule{0pt}{18pt}& & & & \\\midrule
\rule{0pt}{18pt}& & & & \\\midrule
\rule{0pt}{18pt}& & & & \\\midrule
\rule{0pt}{18pt}& & & & \\\midrule
\rule{0pt}{18pt}& & & & \\
\end{longtable}

\subsection{Repeatable review workflow and automation boundary}
\begin{enumerate}
\item \textbf{Register:} freeze the source and render; select a view probe, review audience and relevant evidence records.
\item \textbf{Inspect:} assess every applicable checklist ID, capturing Pass/Fail/N/A and reasons for all exclusions.
\item \textbf{Corroborate:} check material claims against source repositories, infrastructure as code, APIs/events, service catalogs, runbooks or dated telemetry.
\item \textbf{Challenge:} ask a reviewer other than the author to explain system scope, responsibilities, boundaries and principal edge directions.
\item \textbf{Decide:} compute two scores, classify findings, resolve hard gates and document disposition.
\item \textbf{Publish and maintain:} re-render, approve the exact version and repeat the affected tests after material architecture change.
\end{enumerate}
\textbf{Automatable checks:} missing names/IDs, blank labels, diagram compilation, broken links, inconsistent aliases, export text leakage and source/render drift. \textbf{Human judgement:} architectural completeness, correct ownership, trust assumptions, high-risk omissions, semantic accuracy and visual comprehension. CI success never establishes semantic approval.

\subsection{Sign-off record}
\begin{tabularx}{\linewidth}{@{}>{\raggedright\arraybackslash}p{0.42\linewidth}Y@{}}
\toprule\textbf{Field} & \textbf{Reviewer entry}\\\midrule
Diagram ID / view / title & \rule{0pt}{17pt}\\
Scope / environment / current or proposed & \rule{0pt}{17pt}\\
Source revision / rendered artifact & \rule{0pt}{17pt}\\
Applicable / Pass / Fail / N/A & \rule{0pt}{17pt}\\
Checklist coverage / weighted score & \rule{0pt}{17pt}\\
Mandatory gates / open Critical or High & \rule{0pt}{17pt}\\
Decision: Approved / Changes required / Rework & \rule{0pt}{17pt}\\
Reviewer / owner / date / corrective due date & \rule{0pt}{17pt}\\
\bottomrule
\end{tabularx}

\subsection{Source register and companion artifacts}
\begin{itemize}
\item Simon Brown, \emph{Software Architecture Diagram Review Checklist}, official C4 guidance: \url{https://c4model.com/diagrams/checklist}.
\item Simon Brown, \emph{C4 Diagram Notation}, official guidance: \url{https://c4model.com/diagrams/notation}.
\item Simon Brown, \emph{The C4 Model: Visualizing Software Architecture}, O'Reilly, 2026. Taxonomy aligned with supplied appendix text; no facsimile of the book appendix is reproduced.
\item Repository companions: \texttt{scorecards/diagram-quality-scorecard-template.tex} and \texttt{templates/diagram-review-record-template.tex}. Weights and existing bands are preserved while new quality gates are proposed supplements.
\end{itemize}
\end{document}
